<<<<<<< HEAD
\documentclass[landscape,final,a0paper,fontscale=0.29]{baposter}
=======
\documentclass[landscape,final,a0paper,fontscale=0.295]{baposter}
>>>>>>> origin/master

\usepackage[english]{babel}
\usepackage{amsthm}

\usepackage{calc}
\usepackage{graphicx}
\usepackage{amsmath}
\usepackage{amssymb}
\usepackage{relsize}
\usepackage{multirow}
%\usepackage{rotating}
\usepackage{bm}
\usepackage{url}
\usepackage{tikz}
\usetikzlibrary{calc}

%Useful Theorem Environments
\newtheorem{theorem}{Theorem}
\newtheorem{corollary}[theorem]{Corollary}
\newtheorem{lemma}[theorem]{Lemma}
\newtheorem{proposition}[theorem]{Proposition}
\newtheorem{definition}[theorem]{Definition}
\newtheorem{warning}[theorem]{Warning}
\newtheorem{questions}{Question}
\newtheorem{conjecture}[theorem]{Conjecture}
\newtheorem{assumption}{Assumption}
\newtheorem{example}[theorem]{Example}
\newtheorem{construction}[theorem]{Construction}
\newtheorem{quasi-theorem}[theorem]{Quasi-Theorem}

\newtheorem{rem1}[theorem]{Remark}
\newenvironment{remark}{\begin{rem1}\em}{\end{rem1}}

\newtheorem{not1}[theorem]{Notation}
\newenvironment{notation}{\begin{not1}\em}{\end{not1}}

\usepackage{mathptmx} %make math font slightly better

\usepackage{graphicx}
\usepackage{multicol}
\usepackage{palatino}

\newcommand{\captionfont}{\footnotesize}

\graphicspath{{images/}{../images/}}
\usetikzlibrary{calc}

\newcommand{\SET}[1]  {\ensuremath{\mathcal{#1}}}
\newcommand{\MAT}[1]  {\ensuremath{\boldsymbol{#1}}}
\newcommand{\VEC}[1]  {\ensuremath{\boldsymbol{#1}}}
\newcommand{\Video}{\SET{V}}
\newcommand{\video}{\VEC{f}}
\newcommand{\track}{x}
\newcommand{\Track}{\SET T}
\newcommand{\LMs}{\SET L}
\newcommand{\lm}{l}
\newcommand{\PosE}{\SET P}
\newcommand{\posE}{\VEC p}
\newcommand{\negE}{\VEC n}
\newcommand{\NegE}{\SET N}
\newcommand{\Occluded}{\SET O}
\newcommand{\occluded}{o}

\renewcommand{\rmdefault}{ptm} %change all font to Times
\usepackage[english]{babel}

%\rmfamily


\setlength{\columnsep}{1.5em}
\setlength{\columnseprule}{0mm}

%%%%%%%%%%%%%%%%%%%%%%%%%%%%%%%%%%%%%%%%%%%%%%%%%%%%%%%%%%%%%%%%%%%%%%%%%%%%%%%%
% Save space in lists. Use this after the opening of the list
%%%%%%%%%%%%%%%%%%%%%%%%%%%%%%%%%%%%%%%%%%%%%%%%%%%%%%%%%%%%%%%%%%%%%%%%%%%%%%%%
\newcommand{\compresslist}{%
\setlength{\itemsep}{1pt}%
\setlength{\parskip}{0pt}%
\setlength{\parsep}{0pt}%
}

%%%%%%%%%%%%%%%%%%%%%%%%%%%%%%%%%%%%%%%%%%%%%%%%%%%%%%%%%%%%%%%%%%%%%%%%%%%%%%
%%% Begin of Document
%%%%%%%%%%%%%%%%%%%%%%%%%%%%%%%%%%%%%%%%%%%%%%%%%%%%%%%%%%%%%%%%%%%%%%%%%%%%%%

\begin{document}

%%%%%%%%%%%%%%%%%%%%%%%%%%%%%%%%%%%%%%%%%%%%%%%%%%%%%%%%%%%%%%%%%%%%%%%%%%%%%%
%%% Here starts the poster
%%%---------------------------------------------------------------------------
%%% Format it to your taste with the options
%%%%%%%%%%%%%%%%%%%%%%%%%%%%%%%%%%%%%%%%%%%%%%%%%%%%%%%%%%%%%%%%%%%%%%%%%%%%%%
% Define some colors

%\definecolor{lightblue}{cmyk}{0.83,0.24,0,0.12}
%\definecolor{lightblue}{rgb}{0.145,0.6666,1}
\definecolor{lightblue}{rgb}{0.115,0.7266,0.7}

% Draw a video
\newlength{\FSZ}
\newcommand{\drawvideo}[3]{% [0 0.25 0.5 0.75 1 1.25 1.5]
   \noindent\pgfmathsetlength{\FSZ}{\linewidth/#2}
   \begin{tikzpicture}[outer sep=0pt,inner sep=0pt,x=\FSZ,y=\FSZ]
   \draw[color=lightblue!50!black] (0,0) node[outer sep=0pt,inner sep=0pt,text width=\linewidth,minimum height=0] (video) {\noindent#3};
   \path [fill=lightblue!50!black,line width=0pt]
     (video.north west) rectangle ([yshift=\FSZ] video.north east)
    \foreach \x in {1,2,...,#2} {
      {[rounded corners=0.6] ($(video.north west)+(-0.7,0.8)+(\x,0)$) rectangle +(0.4,-0.6)}
    }
;
   \path [fill=lightblue!50!black,line width=0pt]
     ([yshift=-1\FSZ] video.south west) rectangle (video.south east)
<<<<<<< HEAD
=======

>>>>>>> origin/master
    \foreach \x in {1,2,...,#2} {
      {[rounded corners=0.6] ($(video.south west)+(-0.7,-0.2)+(\x,0)$) rectangle +(0.4,-0.6)}
    }
;
   \foreach \x in {1,...,#1} {
     \draw[color=lightblue!50!black] ([xshift=\x\linewidth/#1] video.north west) -- ([xshift=\x\linewidth/#1] video.south west);
   }
   \foreach \x in {0,#1} {
     \draw[color=lightblue!50!black] ([xshift=\x\linewidth/#1,yshift=1\FSZ] video.north west) -- ([xshift=\x\linewidth/#1,yshift=-1\FSZ] video.south west);
   }
   \end{tikzpicture}
}

\hyphenation{resolution occlusions}
%%
\begin{poster}%
  % Poster Options
  {
  % Show grid to help with alignment
  grid=false,
  % Column spacing
  colspacing=1em,
  columns=4,
  % Color style
  bgColorOne=white,
  bgColorTwo=white,
  borderColor=black,
<<<<<<< HEAD
  headerColorOne=black,
  headerColorTwo=lightblue,
=======
  headerColorOne=blue,
  headerColorTwo=blue,
>>>>>>> origin/master
  headerFontColor=white,
  boxColorOne=white,
  boxColorTwo=lightblue,
  % Format of textbox
  textborder=roundedleft,
  % Format of text header
  eyecatcher=true,
  headerborder=closed,
  headerheight=0.1\textheight,
  %textfont=\sc, An example of changing the text font
  headershape=roundedright,
  headershade=shadelr,
  headerfont=\Large\bf\textsc, %Sans Serif
  textfont={\setlength{\parindent}{1.5em}},
  boxshade=plain,
%  background=shade-tb,
  background=plain,
  linewidth=2pt
  }
  % Eye Catcher
    {\includegraphics[height=8em]{images/yale}}
  % Title
  {\bf\textsc{Poster Title}\vspace{0.5em}}
  % Authors
  {\textsc{ Student Names\\ Advised by:}}
  % University logo
  {% The makebox allows the title to flow into the logo, this is a hack because of the L shaped logo.

\setlength{\fboxsep}{0pt}%
\setlength{\fboxrule}{3pt}%
  }

%%%%%%%%%%%%%%%%%%%%%%%%%%%%%%%%%%%%%%%%%%%%%%%%%%%%%%%%%%%%%%%%%%%%%%%%%%%%%%
%%% Now define the boxes that make up the poster
%%%---------------------------------------------------------------------------
%%% Each box has a name and can be placed absolutely or relatively.
%%% The only inconvenience is that you can only specify a relative position
%%% towards an already declared box. So if you have a box attached to the
%%% bottom, one to the top and a third one which should be in between, you
%%% have to specify the top and bottom boxes before you specify the middle
%%% box.
%%%%%%%%%%%%%%%%%%%%%%%%%%%%%%%%%%%%%%%%%%%%%%%%%%%%%%%%%%%%%%%%%%%%%%%%%%%%%%
    %
    % A coloured circle useful as a bullet with an adjustably strong filling
    \newcommand{\colouredcircle}{%
      \tikz{\useasboundingbox (-0.2em,-0.32em) rectangle(0.2em,0.32em); \draw[draw=black,fill=lightblue,line width=0.03em] (0,0) circle(0.18em);}}


%%%%%%%%%%%%%%%%%%%%%%%%%%%%%%%%%%%%%%%%%%%%%%%%%%%%%%%%%%%%%%%%%%%%%%%%%%%%%%
%%%%%%%%%%%%%%%%%%%%%%%%%%%%%%%%%%%%%%%%%%%%%%%%%%%%%%%%%%%%%%%%%%%%%%%%%%%%%%
  %%%%%%%%%%%%%%%%%%%%%%%%%%%%%%%%%%%%%%%%%%%%%%%%%%%%%%%%%%%%%%%%%%%%%%%%%%%%%%

%BACKGROUND SECTION
\headerbox{Background}{name=BG,column=0,row=0,span=4}{
\begin{multicols}{2}
\large
<<<<<<< HEAD
Real-rooted polynomials are these polynomials who have only real roots. We say that a linear transformation $T:\mathbb{R}[x]\to\mathbb{R}[x]$ \emph{preserves real-rootedness}, if $T(f)$ is real-rooted, whenever $f\in\mathbb{R}[x]$ is real-rooted.  Such transformations have been object of many recent research and their theory has had a significant progress in recent years. Notably, Borcea and Br\"{a}nden found their complete characterization. Nevertheless, their result does not shed any light on where on the real line do the roots of $T(f)$ lie as some dependance on the roots of $f$. We analyze this question focusing on two real-rootedness preserves. Most of the results, however, can be extended to a big classes of real-rootedness preservers.
\end{multicols}
}

%%%%%%%%%%%%%%%%%%%%%%%%%%%%%%%%%%%%%%%%%%%%%%%%%%%%%%%%%%%%%%%%%%%%%%%%%%%%%%%

%FIRST MULTICOLUMN
\headerbox{Introduction to making a poster}{name=TL,column=0,span=2,below=BG}{
%%%%%%%%%%%%%%%%%%%%%%%%%%%%%%%%%%%%%%%%%%%%%%%%%%%%%%%%%%%%%%%%%%%%%%%%%%%%%%
\begin{multicols}{2}
\large
\noindent \newtheorem *{question}{Question}
\begin  {question}
The motivating question for my project is: ``can I make a good poster?"
\end {question}

\noindent A poster should not contain too much text -- this is not a good place for your proofs, and pseudocode, it is a gadget that has just enough to point at so you can effectively and crisply explain the essence of your ideas: the problem, some examples, perhaps counterexamples, an approach to a solution, good illustrations.

\begin{definition}
A poster $P$ is said to be \textit{good} if it satisfies all the properties of being good, and is a poster.
\end{definition}

Our first main result is an existence result.
\begin{theorem}
There exist at most finitely many good posters.
\end{theorem}


=======
\noindent Real stable polynomials are polynomials which have only real roots. We say that a linear transformation $T:\mathbb{R}[x]\to\mathbb{R}[x]$ \emph{preserves real stability} if $T(f)$ is real-rooted whenever $f$ is real-rooted.  Such transformations have been object of much recent research and their theory has had significant progress in recent years. Notably, Borcea and Br\"{a}nd\'{e}n found their complete characterization. Nevertheless, their result does not shed any light on where the roots of $T(f)$ lie as a function of the roots of $f$. We analyze this question focusing on two real stability preservers. Most of the results, however, can be extended to a large class of maps.
\end{multicols}
}

%FIRST MULTICOLUMN  
\headerbox{The $1-D$ operator}{name=TL,column=0,span=2,below=BG}{
%%%%%%%%%%%%%%%%%%%%%%%%%%%%%%%%%%%%%%%%%%%%%%%%%%%%%%%%%%%%%%%%%%%%%%%%%%%%%%
\large
\noindent The $1-D$ operator, where $D=\frac{d}{dx}$, is a linear map on the set of real stable polynomials. It was of particular interest in a recent paper solving the Kadison-Singer problem, in part because of its relation to the Laguerre polynomials. We focus primarily on its effect on the differences between the roots of a polynomial.

}

\headerbox{Intuitive picture}{name=TC,column=0,span=2,below=TL}{
%%%%%%%%%%%%%%%%%%%%%%%%%%%%%%%%%%%%%%%%%%%%%%%%%%%%%%%%%%%%%%%%%%%%%%%%%%%%%%
%\begin{multicols}{2}
\large
\noindent Given a polynomial $f$, the roots of ${(1-D)f}$ can be understood as the points of equilibrium in a system of forces resulting from placing a unit charge at each of the roots of $f$ and introducing a constant gravitational force in the $-x$ direction. This is shown in the figure below.
\begin{multicols}{2}
\begin{center}
\includegraphics[width=0.8\columnwidth]{images/charges}
\end{center}
The blue line represents the force at $x$ due to only the blue roots, the red line shows the force due to the red root, and the purple line shows their sum. The red arrow shows how adding the red root repels the old solutions.
>>>>>>> origin/master
\end{multicols}

\noindent This model makes it easy to see the effects of changing the roots of a polynomial, which can be very useful for bounding various quantities. The intuition provided can often be formalized into a more rigorous argument, as was done for Theorem~\ref{thm:rootgaps} below.
%\end{multicols}
}

%%%%%%%%%%%%%%%%%%%%%%%%%%%%%%%%%%%%%%%%%%%%%%%%%%%%%%%%%%%%%%%%%%%%%%%%%%%%%%
%%%%%%%%%%%%%%%%%%%%%%%%%%%%%%%%%%%%%%%%%%%%%%%%%%%%%%%%%%%%%%%%%%%%%%%%%%%%%%
<<<<<<< HEAD
\headerbox{My Generation}{name=BL,column=0,row=3,span=2,below=TL}{
\begin{multicols}{2}
\large
People try to put us d-down (Talkin' 'bout my generation)\\
Just because we get around (Talkin' 'bout my generation)\\
Things they do look awful c-c-cold (Talkin' 'bout my generation)\\
I hope I die before I get old (Talkin' 'bout my generation)\\
\bigskip

This is my generation\\
This is my generation, baby\\
\bigskip

Why don't you all f-fade away (Talkin' 'bout my generation)\\
And don't try to dig what we all s-s-say (Talkin' 'bout my generation)\\
I'm not trying to cause a big s-s-sensation (Talkin' 'bout my generation)\\
I'm just talkin' 'bout my g-g-g-generation (Talkin' 'bout my generation)\\
\bigskip

This is my generation\\
This is my generation, baby\\

\end{multicols}
=======

\headerbox{Asymptotic results}{name=BL,column=0,row=3,span=2,below=TC}{
%\begin{multicols}{1}
\large
\noindent Our main result for this operator is
\begin{theorem}
\label{thm:rootgaps}
	For any real-rooted polynomial $f$, the difference of any two roots of ${(1-D)^n f}$ approaches infinity as $n$ approaches infinity.
\end{theorem}

Interestingly, however, in a relative sense the roots tend to be close together:

\begin{theorem}
\label{thm:rootratio}
	For any polynomial $f$, the ratio of any two roots of ${(1-D)^n f}$ approaches $1$ as $n$ approaches infinity.
\end{theorem}
%\end{multicols}

>>>>>>> origin/master
}

%END OF THIS MULTICOLUMN

%%%%%%%%%%%%%%%%%%%%%%%%%%%%%%%%%%%%%%%%%%%%%%%%%%%%%%%%%%%%%%%%%%%%%%%%%%%%%%
%%%%%%%%%%%%%%%%%%%%%%%%%%%%%%%%%%%%%%%%%%%%%%%%%%%%%%%%%%%%%%%%%%%%%%%%%%%%%%
    \headerbox{The $T_!:x \to x^k/k!$ operator}{name=TR,row=1,column=2,span=2, below=BG}{
%%%%%%%%%%%%%%%%%%%%%%%%%%%%%%%%%%%%%%%%%%%%%%%%%%%%%%%%%%%%%%%%%%%%%%%%%%%%%%
%\begin{multicols}{1}
\large
\noindent
Example:
\[x^4 + 7x^3 + 3x^2 + 2x + 1 \xrightarrow{T_!} \frac{1}{24}x^4 + \frac{7}{6}x^3 + \frac{3}{2}x^2 + 2x + 1\]
\begin{itemize}
<<<<<<< HEAD
 \item Preserves real-rootedness. In fact, if $p(x)$ has real roots $\lambda_1,\cdots,\lambda_n$, then
$T!(p) = [\prod_{\lambda_i \neq 0} (\frac{1}{\lambda_i} - D) \prod_{\lambda_i = 0} D]p(x).$
 \item Links ordinary and exponential functions (used for counting objects in combinatorics):
$\sum_{n\ge 0} a_n x^n \xrightarrow{T_!} \sum_{n\ge 0} \frac{a_n}{n!} x^n$

=======
 \item Preserves real-rootedness. In fact, if $p(x)$ has real roots $\lambda_1,\cdots,\lambda_n$, then we have \\
$T_!(p) = [\prod_{\lambda_i \neq 0} (\frac{1}{\lambda_i} - D) \prod_{\lambda_i = 0} D]p(x).$
 \item Links ordinary and exponential generating functions (used for counting objects in combinatorics):
$\sum_{n\ge 0} a_n x^n \xrightarrow{T_!} \sum_{n\ge 0} \frac{a_n}{n!} x^n$.
  
>>>>>>> origin/master
\end{itemize}
%\end{multicols}
}


%%%%%%%%%%%%%%%%%%%%%%%%%%%%%%%%%%%%%%%%%%%%%%%%%%%%%%%%%%%%%%%%%%%%%%%%%%%%%%
%%%%%%%%%%%%%%%%%%%%%%%%%%%%%%%%%%%%%%%%%%%%%%%%%%%%%%%%%%%%%%%%%%%%%%%%%%%%%%
<<<<<<< HEAD
\headerbox{\vspace{-1cm}Majorization}{name=MJRZ,column=2,row=3,span=2,below=TR}{
\begin{multicols}{2}
\noindent
\large
If the polynomial $p(x)$ has real roots $\lambda_1\leq\cdots\leq\lambda_n$, we associate it with the weakly-increasing vector $\lambda = (\lambda_1,\cdots,\lambda_n)$ of its roots. Suppose $x = (x_1,\cdots, x_n)$ and $y = (y_1,\cdots, y_n)$ are weakly increasing vectors in $\mathbb{R}^n$, then $x$ \emph{majorizes} $y$ (denoted by $x \succ y$) if $\sum_{i=1}^n x_i = \sum_{i=1}^n y_i$ and $\sum_{i=k}^n x_i \geq \sum_{i=}^n y_i$ for $2 \leq k\leq n.$ Given two polynomials $p(x)$ and $q(x)$, we say $p(x)$ \emph{majorizes} $q(x)$, $p \succ q$, if $\lambda \succ \mu$, where $\lambda$ and $\mu$ are the weakly-increasing vectors, associated to $p$ and $q$, respectively.
\begin{theorem}{(Br\"{a}nden)}
=======
\headerbox{Majorization}{name=MJRZ,column=2,row=3,span=2,below=TR}{
\noindent
\large
If a polynomial $p$ has real roots $\lambda_1\leq\cdots\leq~\lambda_n$, we associate it with the weakly-increasing vector $\lambda = (\lambda_1,\cdots,\lambda_n)$ of its roots. Suppose $x = (x_1,\cdots, x_n)$ and $y = (y_1,\cdots, y_n)$ are weakly increasing vectors in $\mathbb{R}^n$, then $x$ \emph{majorizes} $y$ (denoted by $x \succ y$) if $\sum_{i=1}^n x_i = \sum_{i=1}^n y_i$ and $\sum_{i=k}^n x_i \geq \sum_{i=k}^n y_i$ for $2 \leq k\leq n.$ Given two polynomials $p(x)$ and $q(x)$, we say $p$ \emph{majorizes}~$q$, $p \succ q$, if $\lambda \succ \mu$, where $\lambda$ and $\mu$ are the weakly-increasing vectors associated to $p$ and $q$, respectively.
\begin{theorem}{(Br\"{a}nd\'{e}n)}
>>>>>>> origin/master
The linear operator $T_!:~x^k\to\frac{1}{k!}x^k$ preserves majorization, i.e., $T_!(p) \succ T_!(q),$ whenever $p \succ q$.
\end{theorem}
 \begin{corollary}
  $p(x) \succ q(x) \Rightarrow \lambda_{max} T_!(p) \geq \lambda_{max} T_!(q)$.
 \end{corollary}
}

\headerbox{Movement of maximum root}{name=MRT,column=2,row=4,span=2,below=MJRZ}{
\large
\noindent By Vieta's formulas the roots of $p(x) = \sum_{i=0}^n a_i x^i$ have sum $\sum_{i=1}^n \lambda_i = -\frac{a_n}{a_{n-1}}$.
\begin{theorem} For $p(x) = \sum_{i=1} a_i x^i$, we have
\[\lambda_{max}(p(x)) \geq \frac{ \lambda_{max} (T_!(p))}{\lambda_{max}(L_n(x))} \geq -\frac{a_{n-1}}{n a_n},\]
 where $L_n(x)$ is the $n$-th Laguerre polynomial.
\end{theorem}
}

%%%%%%%%%%%%%%%%%%%%%%%%%%%%%%%%%%%%%%%%%%%%%%%%%%%%%%%%%%%%%%%%%%%%%%%%%%%%%%
%%%%%%%%%%%%%%%%%%%%%%%%%%%%%%%%%%%%%%%%%%%%%%%%%%%%%%%%%%%%%%%%%%%%%%%%%%%%%%
\headerbox{References}{name=Ref,column=0,row=5,span=4,below=MRT}{
\noindent

 S. Fisk. {\em Polynomials, roots, and interlacing.} 2008.

 A. Marcus, D. Spielman, and N. Srivastava. {\em Ramanujan graphs and the solution of the Kadison-Singer problem.} Proceedings of ICM, Seoul, 2014.

 J. Borcea and P. Br\"{a}nd\'{e}n. {\em Hyperbolicity preservers and majorization.} 2010.

}


\end{poster}

\end{document}
